\section{The Mukai pairing and adjoint kernels}
\label{sec:adjoint}

In this section we define the Mukai pairing on the Hochschild homology
of a space and show that the push-forwards of adjoint kernels are
themselves adjoint linear maps with respect to this pairing.

%\subsection{Definition of the Mukai pairing.}
First observe from Section~\ref{subsec:adj2functors} that we have two
isomorphisms:
\[\tau_R,\tau_L\colon\HH_\blob(X) = \Hom^{-\blob}(\SerK_X^{-1}, \Id_X) \xrightarrow{\sim}\Hom^\blob(\Id_X, \SerK_X),\] 
given by
\[\tau_R\left(\pic{vinHH}\right):=\pic{tauRv}\qquad\text{and} \quad
\tau_L\left(\pic{vprimeinHH}\right):=\pic{tauLvprime}.\]
Note that this differs slightly from the given definition, but we have used the uniqueness of adjoints.  The above isomorphisms allow the definition of the Mukai pairing as
follows.

\begin{Definition} %\textbf{(Mukai pairing).}
  The \emph{Mukai pairing} on the Hochschild homology of a space $X$ is the
  map
  \[ \MP{{-}}{{-}} \colon \HH_\blob(X) \otimes \HH_\blob(X) \ra \k, \]
  defined by
  \[ \MP{v}{v'}:= \Tr\left(\tau_R(v) \circ \tau_L(v')\right). \]
  Diagrammatically, this is
  \[\MP{\pic{vinHH}}{\pic{vprimeinHH}}:=
  \Tr\left(\pic{tauRv}\pic{tauLvprime}\right).\]
\end{Definition}
\medskip

\noindent
Observe that as $\tau_R$ and $\tau_L$ are both isomorphisms and as the
Serre pairing is non-degenerate, it follows that the Mukai pairing is
non-degenerate.

%\subsection{Adjoint kernels and adjoint maps.}
We can now easily show that adjoint kernels give rise to adjoint maps
between the corresponding Hoch\-schild homology groups.

\begin{Theorem}[Adjointness]
\label{thm:adjointness}
  If $\Phi\colon X\to Y$ and $\Psi\colon Y\to X$ are adjoint kernels,
  $\Psi\dashv\Phi$, then the corresponding push forwards are adjoint
  with respect to the Mukai pairing in the sense that for all $w\in
  \HH_\blob(X)$ and $v\in \HH_\blob(Y)$ we have
  \[\MP{\Psi_*(v)}{w}=\MP{v}{\Phi_*(w)}.\]
\end{Theorem}

\begin{Proof}
  Note first that $\Psi_*=\Phi^*$, by Theorem~\ref{thm:adjadj}.  Thus
  \begin{align*}
    \MP{\Psi_*(v)}{w}&=\MP{\Phi^*(v)}{w}
    =\Tr\left(\pic{tauRPhiv}\pic{tauLw}\right)\\
    &=\Tr\left(\pic{tauRv}\pic{IdPhiarrow}\pic{tauLw}\right)
    =\Tr\left(\pic{tauRv}\pic{tauLPhiw}\right)\\
    &=\MP{v}{\Phi_*(w)}.\qedhere
  \end{align*}
\end{Proof}

\begin{Corollary}
  If the integral kernel $\Phi\colon X \ra Y$ induces an equivalence
  on derived categories, then $\Phi_*\colon\HH_\blob(X) \ra
  \HH_\blob(Y)$ is an isometry.
\end{Corollary}

\begin{Proof}
  If $\Phi$ induces an equivalence, then it has a left adjoint
  $\Psi\colon Y\to X$ which induces the inverse, so
  $\Psi\circ\Phi\cong \Id_X$, and we know that $(\Id_X)_*$
  is the identity map.  Thus
  \[ \MP{\Phi_* v}{\Phi_* w} = \MP{\Psi_* \Phi_* v}{w} =
  \MP{(\Psi\circ \Phi)_* v}{w}= \MP{v}{w}.\qedhere \]
\end{Proof}
